\section{三、模型假设}\label{ux4e09ux6a21ux578bux5047ux8bbe}

为使复杂的工业生产调度问题具有清晰严谨的数学结构，结合实际运作场景，提出以下合理假设：

\begin{enumerate}
\def\labelenumi{\arabic{enumi}.}
\item
  \textbf{周期离散化与时序同步假设}：
  以周为基本离散决策时间单位。每周初集中下达当周物料生产计划，周末统一盘点当周结算库存量与缺货量。各周的生产加工在周内平稳进行。
  \emph{理由}：符合制造业周度主生产计划（MPS）运作规律及竞赛附录明确规定。
\item
  \textbf{提前期确定性交付假设}：
  在问题2和问题3中，物料加工与物流提前期恒定为 \(L=1\) 周，即第 \(t\)
  周初下达的生产计划 \(P_t\)，在第 \(t+1\)
  周初准时送达并可供使用；在问题4中，提前期为确定的 \(L=k\)
  周。生产过程无废品与交货延误。
  \emph{理由}：符合题目设定，聚焦于需求不确定性下的调度优化主干。
\item
  \textbf{信息因果律与决策不可预知性假设}： 在第 \(t\) 周初制定生产计划
  \(P_t\) 时，决策系统严格只能获取第 \(t-1\)
  周及以前的所有历史需求、库存和预测数据，严禁使用第 \(t\)
  周及以后的任何实际需求。
  \emph{理由}：严格遵循附录(2)红线，杜绝任何未来信息泄漏。
\item
  \textbf{丢失销售（Lost Sales）与当周独立结算假设}：
  当周物资供给无法满足的客户需求计入当周缺货量
  \(B_t\)；由于电子消费品快速迭代，未满足的需求不积压结转至下周（即采用丢失销售机制），各周服务水平根据当周实际需求独立核算。
  \emph{理由}：符合附录(3)(4)公式定义：\(\text{服务水平} = 1 - \frac{\text{缺货量}}{\text{实际需求量}}\)，且符合现货满足率（Fill
  Rate）评价准则。
\item
  \textbf{产线独立排产假设}： 筛选出的 6
  种重点物料在各自对应的产线或工位独立安排生产，不存在跨物料的产能瓶颈互斥冲突。
  \emph{理由}：符合题目对各物料分别制定生产计划与填报表格的要求。
\end{enumerate}
